\section{Round-four referee closure: the complete vector--roof packet}
\label{sec:r3-a2}

This section replaces the imported spectral packet and the unproved arithmetic
paragraph in the body.  The proof is organized as a finite Banach bundle over
the radius interval, an explicit temporal non-integrability rectangle, and a
three-frequency Fourier argument.  The constants are uniform for
\[
 \mathcal I=[9/20,47/100].
\]

\subsection{Uniform geometry and the moving-cut bundle}

The triangular lattice has fundamental area \(A_\Lambda=\sqrt3/2\).  For a primitive lattice direction \(v\), the distance between adjacent lattice rows
parallel to \(v\) is
\[
 d_\perp(v)=\frac{A_\Lambda}{|v|}
 \le\frac{\sqrt3}{2}<2R
\]
for every \(R\in[9/20,47/100]\).  Thus every rational direction is blocked;
an irrational direction cannot support a periodic open strip.  Also
\[
 g_*=1-2(47/100)=3/50>0,
 \qquad
 2(9/20)-\sqrt3/2>0,
\]
so the tables remain uniformly separated from overlap and corridor
thresholds.  Curvature lies in \([100/47,20/9]\).

\begin{lemma}[Uniform finite-horizon and finite-iterate expansion geometry]
\label{lem:r3-a2-geometry}
There are
\[
 0<\tau_-<\tau_+<\infty,
 \quad \Lambda_*>1,
 \quad C_{\rm d}<\infty,
 \quad n_*\in\mathbb N,
 \quad 0<\vartheta_*<1,
\]
which work for every \(R\in\mathcal I\).  Every regular free flight lies in
\([\tau_-,\tau_+]\); the standard stable and unstable cones are invariant;
homogeneous unstable curves have uniform expansion and distortion; and the
\(n_*\)-step expansion sum is at most \(\vartheta_*\).
\end{lemma}

\begin{proof}
The lower flight bound follows from \(g_*>0\).  If no common upper bound
existed, a sequence of increasingly long free segments would have a limiting
base point and direction.  Its lift would give an open Euclidean strip free of
lattice disks.  A strip direction containing two lattice points has width at
most \(d_\perp(v)\le\sqrt3/2<2R\), while an irrational strip has lattice
translates dense in the transverse torus and hence cannot remain open.
Therefore every table has finite horizon.  Compactness of the radius interval,
base point, and direction gives one \(\tau_+\).

The mirror equation and the curvature/flight bounds give common invariant
cones.  On the homogeneity strip
\(k^{-2}\le\cos\varphi<(k-1)^{-2}\), the first two derivative and distortion
bounds have the standard polynomial weights in \(k\), with coefficients
depending only on \(\tau_\pm\) and the curvature interval.  For each radius,
the growth lemma gives an iterate \(n(R)\) whose inverse-expansion sum is
strictly below one.  The same inequality persists on a radius neighborhood by
the moving singularity charts.  A finite subcover and a common multiple of
the finitely many iterates give \(n_*\) and \(\vartheta_*<1\).  Renorming by
\[
 \|h\|_*=
 \sum_{j=0}^{n_*-1}\rho^{-j}\|\mathcal L_R^jh\|
\]
converts this finite-iterate contraction into the uniform one-step
Lasota--Yorke form used below.
\end{proof}

Use common angular coordinates
\((r,\varphi)\in\mathbb T\times(-\pi/2,\pi/2)\).  For \(n\ge0\), let
\(\mathscr S_{R,n}\) be the collection of connected components of
\(T_R^{-n}\mathscr S_{R,0}\), cut at intersections with all lower-depth
components and with homogeneity boundaries.  A component
\(\gamma\in\mathscr S_{R,n}\) carries its depth \(d(\gamma)=n\), its two
one-sided traces, and its orientation induced by the collision coordinates.

For an admissible stable curve \(W\), write
\[
 \langle h,\psi\rangle_W=\int_Wh\psi\,dm_W.
\]
The ordinary strong norm contains the standard stable, unstable-holonomy, and
short-curve terms.  Moving a singularity creates an additional trace current,
so define
\[
 \mathcal C_R
 =\left\{J=(J_\gamma)_{\gamma}:
 \|J\|_{\mathcal C_R}
 =\sum_{\gamma}\eta^{d(\gamma)}
 (1+d(\gamma))^{-3}
 \|J_\gamma\|_{(C^{1+\alpha}(\gamma))^*}<\infty
 \right\},
\]
where \(\eta>1\) is chosen below, and put
\[
 \mathbb B_R=\mathcal B_R\oplus\mathcal C_R,
 \qquad
 \mathbb B_{R,w}=\mathcal B_{R,w}\oplus\mathcal C_{R,w}.
\]
A physical density is embedded as \((h,0)\); a derivative in \(R\) generally
has a nonzero current component.

\begin{lemma}[Current-augmented moving-cut bundle]
\label{lem:r3-a2-atlas}
There are \(1<\eta<\Lambda_*\), a finite cover \((U_j)\) of
\(\mathcal I\), fixed reference spaces
\(\mathbb B_j\hookrightarrow\mathbb B_{j,w}\), and identifications
\[
 \mathbb T_{R,j}:\mathbb B_R\longrightarrow\mathbb B_j,
 \qquad R\in U_j,
\]
with uniformly bounded inverses, such that:
\begin{enumerate}[label=(\roman*)]
\item every iterated singularity component is continued until its first
      tangency or intersection, and a birth/death at that radius is represented
      by an oriented current rather than by a discontinuous change of norm;
\item the charted transfer operators
\(\mathbb L_R=\mathbb T_{R,j}\mathcal L_R\mathbb T_{R,j}^{-1}\) are
\(C^1\) from \(U_j\) to
\(\mathcal L(\mathbb B_j,\mathbb B_{j,w})\);
\item for common \(C,\rho<1\),
\[
 \|\mathbb L_R^nh\|_{\mathbb B_R}
 \le C\rho^n\|h\|_{\mathbb B_R}
     +C\|h\|_{\mathbb B_{R,w}},
 \qquad
 \|\mathbb L_R^nh\|_{\mathbb B_{R,w}}
 \le C\|h\|_{\mathbb B_{R,w}};
\]
\item the strong unit ball is relatively compact in the weak space, uniformly
      in \(R\).
\end{enumerate}
The derivative is the sum of the ordinary branch derivative and the absolutely
convergent moving-boundary current series.
\end{lemma}

\begin{proof}
Fix a finite depth \(N\).  Away from finitely many parameter values at which a
tangency or multiple intersection occurs, all components of
\(\bigcup_{n\le N}\mathscr S_{R,n}\) continue by the implicit-function
theorem with uniform \(C^2\) bounds in homogeneous coordinates.  At a critical
radius the difference between the two one-sided pullbacks is, by the transport
formula for a moving cut,
\[
 \frac d{dR}\int_{W_R}h\psi
 =\int_{W_R}\partial_R(h\psi)
  +\sum_{x\in\partial W_R}
  \dot x_R\,[h\psi]_x .
\]
The endpoint terms are precisely elements of \(\mathcal C_R\).  Thus births,
deaths, and crossings are continuous in the augmented weak norm.

For arbitrary depth, the number of children and their inverse unstable
Jacobians satisfy the finite-iterate growth inequality from
Lemma~\ref{lem:r3-a2-geometry}.  Choosing
\(1<\eta<\Lambda_*\) sufficiently close to one gives
\[
 \sup_W\sum_{V\subset T_R^{-n_*}W}
 \eta^{d(V)-d(W)}(J^uT_R^{n_*}|_V)^{-1}
 \le\vartheta_1<1.
\]
The polynomial depth factor absorbs intersections of equal depth.  Therefore
the trace-current series created by differentiating all preimages is
absolutely summable and its tail after depth \(N\) is \(O(\vartheta_1^{N/n_*})\).
This proves the \(C^1\) operator statement after passing from the finite-depth
charts to their limit.

The usual stable/unstable Lasota--Yorke estimates control the distribution
part.  The displayed weighted expansion controls the current part, while the
current-to-density off-diagonal block is weak because it is supported on cut
curves.  Taking a common multiple of the finite-iterate estimates gives (iii).
Arzel\`a--Ascoli on stable test functions, diagonal compactness in depth, and
the geometric current tail give (iv).  A finite parameter subcover completes
the construction.
\end{proof}

\subsection{Uniform twisted spectrum}

The displacement \(\kappa_R\) is constant on every continuity component and
has finite range.  On a homogeneous branch, the roof \(\tau_R\) and its first
two stable/unstable derivatives obey the same strip weights as the Jacobian.

\begin{lemma}[Vector--roof multiplier bound]
\label{lem:r3-a2-multiplier}
There are \(\delta>0\) and \(C<\infty\) such that, for complex
\((\xi,s)\in\mathbb C^2\times\mathbb C\) with
\(|\xi|+|s|<\delta\), multiplication by
\[
 e^{\langle\xi,\kappa_R\rangle-s\tau_R}
\]
is analytic from the strong space to itself and satisfies
\[
 \|M_{R,\xi,s}-I\|_{\mathcal B_R\to\mathcal B_R}
 +\|M_{R,\xi,s}-I\|_{\mathcal B_{R,w}\to\mathcal B_{R,w}}
 \le C(|\xi|+|s|)
\]
uniformly in \(R\).
\end{lemma}

\begin{proof}
The displacement factor is branchwise constant.  For the roof factor, expand
the exponential in its absolutely convergent power series.  The product
estimate in the stable norm uses bounded oscillation on a homogeneous stable
curve; the unstable norm uses the holonomy estimate for the roof.  The strip
weights from Lemma~\ref{lem:r3-a2-geometry} make the sums uniform.  Termwise
differentiation proves analyticity and the displayed Lipschitz estimate.
\end{proof}

\begin{theorem}[Uniform complex spectral bundle]
\label{thm:r3-a2-spectrum}
For \(|\xi|+|s|<\delta\),
\[
 \mathcal L_{R,\xi,s}h
 =\mathcal L_R\left(
 e^{\langle\xi,\kappa_R\rangle-s\tau_R}h\right)
\]
has a simple leading eigenvalue \(e^{P_R(\xi,s)}\), and
\[
 \mathcal L_{R,\xi,s}^n
 =e^{nP_R(\xi,s)}\Pi_{R,\xi,s}+N_{R,\xi,s}^n,
 \qquad
 \|N_{R,\xi,s}^n\|_{\mathcal B_R\to\mathcal B_{R,w}}
 \le C\rho^n e^{n\Re P_R(\xi,s)}
\]
with one \(\rho<1\).  The pressure and projector are analytic in
\((\xi,s)\), continuously differentiable in \(R\), and agree on chart
overlaps.
\end{theorem}

\begin{proof}
The growth lemma and distortion bounds give, for suitable \(b\),
\[
 \|\mathcal L_R^nh\|_{\mathcal B_R}
 \le C\rho_0^n\|h\|_{\mathcal B_R}+C\|h\|_{\mathcal B_{R,w}},
 \qquad
 \|\mathcal L_R^nh\|_{\mathcal B_{R,w}}
 \le C\|h\|_{\mathcal B_{R,w}},
\]
with uniform constants.  Compactness of the strong unit ball in the weak
space gives uniform quasicompactness.  Mixing of the finite-horizon collision
map makes the eigenvalue one simple.  Lemma~\ref{lem:r3-a2-multiplier} and a
resolvent contour lying in the common spectral gap give analytic perturbation
with a common radius.  The moving-cut atlas identifies the resolvent pairings
on overlaps, hence also the Riesz projectors.  Differentiating the charted
operators, including the boundary-motion term already represented by the cut
norm, gives continuous \(R\)-dependence.
\end{proof}

The same proof applies to a fixed finite collision cylinder inserted in the
middle of the orbit.  In particular left and right eigenvectors can be
paired with a central-window observable without replacing it by an evolved
smooth observable.

\subsection{Certified temporal non-integrability on the induced quotient}

The high-frequency argument is carried out on the stable quotient of a
finite-connector Young magnet.  This point is essential: the billiard map is
invertible, whereas the quotient return map is a uniformly expanding map with
countably many inverse branches.  The spectral data of the quotient and of
the anisotropic collision operator agree on the leading eigenspace.

Let
\[
 c_0=(0,0),\qquad c_1=(1,0),\qquad
 c_2=(1/2,\sqrt3/2),
\]
and let \(c=(1/2,\sqrt3/6)\) be the center of the triangular gap.  The three
inward normal collision points are
\[
 p_0=(\sqrt3R/2,R/2),\quad
 p_1=(1-\sqrt3R/2,R/2),\quad
 p_2=(1/2,\sqrt3/2-R).
\]
They form an equilateral triangle of side
\[
 d_R=1-\sqrt3R\in[1-47\sqrt3/100,1-9\sqrt3/20],
\]
which is uniformly positive.  The two oriented words
\[
 w_+=(0,1,2,0),\qquad w_-=(0,2,1,0)
\]
are regular specular return polygons.  At every internal vertex the two
adjacent unit directions have the inward disk normal as their angle bisector.

For endpoints with arclength coordinates \((r,r')\) on a common small arc
about \(p_0\), let \(\mathscr L_{R,\pm}(r,r')\) be the total Euclidean length
of the \(w_\pm\) polygon after the two internal collision coordinates have
been eliminated by the specular critical equations.

\begin{lemma}[Uniform polygon generating functions]
\label{lem:r3-a2-polygon}
There are a radius-independent endpoint interval \(J\), a neighborhood of the
symmetric collision polygon, and constants \(c_H,C_H>0\) such that the
internal critical points defining \(\mathscr L_{R,\pm}\) exist uniquely and
are \(C^3\) in \((R,r,r')\).  The internal Hessian of the polygonal length
satisfies
\[
 c_H I\le D^2_{\rm int}\mathscr L_{R,\pm}\le C_H I.
\]
At the symmetric diagonal endpoint,
\[
 \partial_{r'}\mathscr L_{R,+}(0,0)
 -\partial_{r'}\mathscr L_{R,-}(0,0)=1.
\]
Consequently, after shrinking \(J\) once,
\[
 \inf_{R\in\mathcal I}\inf_{(r,r')\in J^2}
 \left|\partial_{r'}(\mathscr L_{R,+}-
 \mathscr L_{R,-})(r,r')\right|\ge\frac34.
\]
\end{lemma}

\begin{proof}
For a segment joining two strictly convex boundary components, the mixed
second derivative of its length generating function is
\[
 -\frac{\cos\varphi\cos\varphi'}{\tau},
\]
and each diagonal second derivative is the sum of the adjacent free-flight
term and the curvature term \(R^{-1}\cos^{-1}\varphi\).  On the two triangular
polygons all flight lengths lie in a compact subinterval of \((0,\infty)\)
and all collision angles stay a fixed distance from grazing, uniformly for
\(R\in\mathcal I\).  Strict dispersing curvature therefore makes the two by
two internal Hessian uniformly positive.  The implicit-function theorem gives
the stated critical points and regularity.

After the internal variables are eliminated, the envelope theorem says that
the terminal derivative is the tangential component of the incoming unit
momentum.  At \(p_0\) the positively oriented last leg has direction
\((-1/2,-\sqrt3/2)\), the negatively oriented last leg has direction
\((-1,0)\), and the positive unit tangent to the disk is
\((-1/2,\sqrt3/2)\).  Their tangential components are respectively
\(-1/2\) and \(1/2\); reversing the endpoint convention changes both signs
and leaves the absolute difference equal to one.  Uniform continuity on the
compact radius interval gives the \(3/4\) bound.
\end{proof}

Choose one common regular connector word \(c_*\) from the terminal arc back to
the base of the magnet.  Precompose both triangular words with \(m\) copies of
\(c_*\), where \(m\) is fixed below.  Let \(h_{R,+},h_{R,-}\) be the resulting
inverse branches of the stable-quotient return map and let
\(\mathcal R_{R,+},\mathcal R_{R,-}\) be their return roofs.  Stable holonomy
identifies their terminal endpoints.

\begin{lemma}[Uniform temporal-distance derivative]
\label{lem:r3-a2-uni}
There are a fixed connector length \(m\), a common quotient interval \(J_0\),
and \(c_{\rm UNI}>0\) such that the temporal distance
\[
 \Psi_R(x)=\mathcal R_{R,+}(h_{R,+}x)
 -\mathcal R_{R,-}(h_{R,-}x)
 +\sum_{n\ge1}
 \left(\tau_R(F_R^{-n}h_{R,+}x)-
       \tau_R(F_R^{-n}h_{R,-}x)\right)
\]
is \(C^1\) on \(J_0\) and satisfies
\[
 \inf_{R\in\mathcal I}\inf_{x\in J_0}|\Psi_R'(x)|
 \ge c_{\rm UNI}.
\]
One may take \(c_{\rm UNI}=1/4\).
\end{lemma}

\begin{proof}
The first difference is the reduced polygonal generating-function difference
of Lemma~\ref{lem:r3-a2-polygon}, transported by stable holonomy.  The
holonomy Jacobian differs from one by a uniformly bounded factor and can be
made arbitrarily close to one by shrinking the common interval.  The common
connector makes the derivative of the remaining past tail at most
\[
 C\Lambda_*^{-m}\sum_{n\ge0}\Lambda_*^{-n}.
\]
Choose \(m\) so this is at most \(1/4\), and then shrink the interval so the
holonomy error is at most \(1/4\).  The \(3/4\) polygon derivative then leaves
a lower bound \(1/4\).  All choices are uniform because the geometry and
hyperbolicity constants are uniform in \(R\).
\end{proof}

\subsection{Frequency-adapted Dolgopyat contraction}

Let \(b=1+|t|\).  On the quotient return map and on the augmented collision
space use the norm
\[
 \|h\|_{b}
 =\|h\|_{w}+b^{-1}\|h\|_{s}
   +b^{-1/2}\|h\|_{u}
   +b^{-1}\|h\|_{\mathcal C_R}.
\]
The current factor is scaled by \(b^{-1}\) because differentiating a phase
across a moving cut produces one factor \(b\), exactly compensated by this
norm.

\begin{lemma}[Uniform return and cancellation block]
\label{lem:r4-a2-dolgopyat-block}
There are \(C_0,c_0>0\) such that, for every \(b\ge2\), every normalized
standard family has, within
\(N_b\le C_0\log b\) iterates, a subfamily of mass at least \(c_0\) returning
to the common UNI interval of Lemma~\ref{lem:r3-a2-uni}.  On that subfamily the
two inverse branches may be paired so that their twisted sum loses a fixed
fraction \(c_0\) in the \(b\)-norm.  All constants are uniform in
\(R\in\mathcal I\) and \(u\in[-\pi,\pi]^2\).
\end{lemma}

\begin{proof}
The finite-iterate growth lemma first expands every short unstable curve to a
proper standard family.  Exponential tails of the common magnet and the finite
connector word give a return to the UNI base with probability bounded below
within \(C_0\log b\) iterates; the logarithmic delay is the time needed for an
initial curve of length \(b^{-1}\) to reach the proper scale.

On the UNI interval the phase difference of the two branches is
\(t\Psi_R+u\cdot\Delta\kappa\).  The lattice term is locally constant and
\(|\partial_x(t\Psi_R)|\ge c_{\rm UNI}|t|\).  Partition into intervals of
length between \(c/b\) and \(C/b\).  On alternate subintervals the two unit
complex phases have angular separation bounded away from zero.  The branch
Jacobians and conditional densities have bounded distortion, so the elementary
inequality
\[
 |z_1+z_2|\le(1-c_0)(|z_1|+|z_2|)
\]
applies after a \(C^1\) cutoff whose derivative is \(O(b)\).  The derivative
cost is absorbed by \(b^{-1}\|\cdot\|_s\).  The one-sided cut traces obey the
same pairing; their derivative cost is absorbed by the current term.  The
unpaired boundary intervals have total relative mass \(O(b^{-1})\).  This
proves the fixed loss on a mass \(c_0\) subfamily.
\end{proof}

\begin{theorem}[Joint aperiodicity and high-frequency contraction]
\label{thm:r3-a2-high}
If \((u,t)\in[-\pi,\pi]^2\times\mathbb R\) and
\(\mathcal L_{R,iu,it}\) has peripheral spectral radius one, then
\((u,t)=(0,0)\).  Moreover, for \(|t|\ge1\),
\[
 \|\mathcal L_{R,iu,it}^n\|_{\mathbb B_R\to\mathbb B_{R,w}}
 \le C(1+|t|)^A
 \exp\left\{-\frac{cn}{\log(2+|t|)}\right\},
\]
uniformly in \(R\) and \(u\).
\end{theorem}

\begin{proof}
A peripheral eigenvector has constant modulus on the mixing stable quotient.
Periodic branch identities imply that
\(u\cdot S_n\kappa_R-tS_n\tau_R\) is constant modulo \(2\pi\) on every full
branch.  Differentiating the two polygon branches in
Lemma~\ref{lem:r3-a2-uni} forces \(t=0\); the four open displacement branches
then force \(u\in2\pi\mathbb Z^2\), hence \(u=0\) in the fundamental domain.

For the quantitative estimate, apply
Lemma~\ref{lem:r4-a2-dolgopyat-block} every \(N_b\) iterates.  The uncancelled
mass is made proper again by the growth lemma, while the augmented
Lasota--Yorke estimate controls the strong and current components.  Thus
\[
 \|\mathcal L_{R,iu,it}^{kN_b}h\|_b
 \le C(1-c_1)^k\|h\|_b+C b^A\rho^{kN_b}\|h\|_b.
\]
Writing \(k=\lfloor n/N_b\rfloor\) gives the displayed bound.  Compact
nonzero frequency annuli have spectral radius uniformly below one by the
common augmented spectral gap and the absence of peripheral modes.
\end{proof}

\subsection{Joint lattice--nonlattice local limit theorem}

For a real simple source \((\xi,s)\), define
\[
 a_{R,\xi,s}=\partial_\xi P_R(\xi,s),
 \qquad
 \bar\tau_{R,\xi,s}=-\partial_sP_R(\xi,s)>0.
\]
Let \(J=\operatorname{diag}(1,1,-1)\).  The Hessian of \(P_R\) in
\((\xi,s)\) is \(J\Sigma J\), where \(\Sigma\) is the asymptotic covariance
of \((\kappa_R,\tau_R)\) under the tilted phase.

\begin{lemma}[Uniform Liv\v{s}ic alternative and covariance]
\label{lem:r3-a2-cov}
On every compact real simple-source set,
\[
 cI_3\le\Sigma_{R,\xi,s}\le CI_3
\]
with constants independent of \(R\).  If a real dynamically H\"older
observable on the collision map has zero asymptotic variance, then its lift to
the common induced quotient is a H\"older coboundary; conversely every such
coboundary has zero variance.
\end{lemma}

\begin{proof}
The augmented spectral theorem gives a simple analytic eigenvalue for the
one-parameter real twist by any linear combination
\(g=u\cdot\kappa_R-t\tau_R\).  If the second derivative at zero vanishes,
the normalized twisted eigenvalue has modulus one to second order.  Positivity
of the spectral measure and the spectral gap then imply that the martingale
part in the Gordin decomposition vanishes.  Summing the exponentially
convergent reduced-resolvent series gives
\(g-\nu(g)=H-H\circ T_R\) with \(H\) dynamically H\"older on every regular
component and matching one-sided traces in the current completion.

Lifting to the finite-connector quotient preserves periodic sums.  The two UNI
periodic words force \(t=0\); the four primitive displacement words force
\(u=0\).  Thus no nonzero direction has zero variance.  Continuity of the
Hessian in \((R,\xi,s)\) and compactness yield the uniform lower bound; the
upper bound follows from exponential correlation summability.
\end{proof}

\begin{theorem}[Uniform vector--clock local limit]
\label{thm:r3-a2-llt}
Let \(F\) be a bounded central cylinder.  Uniformly for \(R\), for real tilts
in a compact simple branch, and for
\[
 k_n-na_{R,\xi,s}=O(\sqrt n),
 \qquad
 t_n-n\bar\tau_{R,\xi,s}=O(\sqrt n),
\]
let \(J_n=[t_n-b_n,t_n+b_n]\), where
\(b_n\to\infty\) and \(b_n=o(n)\).  Then
\[
 \begin{aligned}
 &\mathbb E_{R,0}
 \left[F;\ S_n\kappa_R=k_n,\ S_n\tau_R\in J_n\right]\\
 &\quad=e^{-nI_R(k_n/n,t_n/n)}
 \frac{2b_n}{(2\pi n)^{3/2}\sqrt{\det\Sigma_{R,\xi,s}}}
 \left(\nu_{R,\xi,s}(F)+o(1)\right).
 \end{aligned}
\]
The estimate is uniform for \(F\) in the unit ball of a fixed cylinder space
and includes \(b_n=o(\sqrt n)\).
\end{theorem}

\begin{proof}
Tilt to the real saddle matching \((k_n/n,t_n/n)\).  Fourier inversion uses
\((u,t)\in[-\pi,\pi]^2\times\mathbb R\).  In the central region
\(|(u,t)|\le n^{-2/5}\), analytic perturbation on the augmented bundle gives
\[
 P_R(\xi+iu,s+it)-P_R(\xi,s)
 =iu\cdot a-it\bar\tau
 -\frac12\langle(u,-t),\Sigma(u,-t)\rangle
 +O(|(u,t)|^3),
\]
with the same expansion for the two projectors surrounding the central
insertion.  Scaling by \(n^{-1/2}\) gives the three-dimensional Gaussian
integral.

Choose a fixed \(C^\infty\) mollifier of width one and upper/lower smooth
approximations of \(1_{J_n}\) which differ only in two boundary strips of
width two.  The central Gaussian estimate bounds the boundary contribution by
\(O(n^{-1/2})\), whereas the main roof-window mass is
\(2b_n(2\pi n)^{-1/2}\); the relative error is therefore \(O(b_n^{-1})\).
The Fourier transforms of the smooth approximants decay faster than any
power with constants independent of \(b_n\).

On
\(n^{-2/5}<|(u,t)|\le\delta\), the covariance lower bound gives
\(e^{-cn|(u,t)|^2}\).  On compact minor arcs, the uniform spectral radius is
below one.  On \(|t|\ge1\), Theorem~\ref{thm:r3-a2-high} together with the
rapid Fourier decay makes the integral
\(O(n^{-M})\) after splitting into blocks of length \(C\log(2+|t|)\).
These estimates are \(o(b_nn^{-3/2})\) for every sequence in the theorem.
Undoing the real tilt supplies the rate exponent and proves the formula.
\end{proof}

\begin{corollary}[Two-sided current--clock conditioning]
\label{cor:r3-a2-conditioning}
Under the hypotheses of Theorem~\ref{thm:r3-a2-llt}, the conditional law of
every fixed central window given
\(S_n\kappa_R=k_n\) and \(S_n\tau_R\in J_n\) converges in total variation to
\(\nu_{R,\xi,s}\).  Uniformly on the declared compact sets, its error is
\[
 O\left(n^{-1/2}+b_n/n+b_n^{-1}\right)
\]
plus the error in matching the prescribed means to the real saddle.
\end{corollary}

\begin{proof}
Apply the local limit theorem to the numerator with a central cylinder and to
the denominator with \(F=1\).  The fixed-width smoothing gives the
\(b_n^{-1}\) term, the cubic central expansion gives \(n^{-1/2}\), and
variation of the saddle over an interval of relative width \(b_n/n\) gives
the remaining term.  The denominator has a uniform positive Gaussian main
term, so division is stable.
\end{proof}

\subsection{Physical pressure root}

Because \(\partial_sP_R(0,0)=-\mu_R(\tau_R)<0\), the equation
\[
 P_R(\xi,\mathcal H_R(\xi))=0
\]
has a uniform analytic solution near zero.

\begin{theorem}[Closed homological pressure theorem]
\label{thm:r3-a2-main}
The pressure root is analytic jointly in the homological field and
continuously differentiable in \(R\), and
\[
 \nabla\mathcal H_R(\xi)
 =\frac{\nu_{R,\xi}(\kappa_R)}{\nu_{R,\xi}(\tau_R)},
\]
\[
 \nabla^2\mathcal H_R(\xi)
 =\frac{\operatorname{AVar}_{\nu_{R,\xi}}
 (\kappa_R-\nabla\mathcal H_R(\xi)\tau_R)}
 {\nu_{R,\xi}(\tau_R)}.
\]
The Hessian is uniformly positive definite on a common field ball.  Exact
vector-current and submacroscopic clock conditioning on a central window are
given by Corollary~\ref{cor:r3-a2-conditioning}.
\end{theorem}

\begin{proof}
The implicit-function theorem applies to
Theorem~\ref{thm:r3-a2-spectrum}.  Differentiating the root once and twice
gives the two formulas; the second is the Schur complement of the joint
vector--roof covariance.  Lemma~\ref{lem:r3-a2-cov} and the positive mean roof
give uniform positivity.  The conditioning statement was proved above.
\end{proof}
